\documentclass[10pt,a4paper]{amsart}
\usepackage[margin=19mm]{geometry}
\usepackage{amsmath,amssymb,mathtools}
\usepackage[scr=rsfs,cal=euler]{mathalfa}
\usepackage{xfrac}
\usepackage[unicode,hidelinks,bookmarks=false]{hyperref}
\newcommand{\ie}{i.e.\ }
\newcommand{\cf}{cf.\ }
\newcommand{\resp}{respectively\ }
\newcommand{\Subsection}{\subsection}
\providecommand{\nobreakdash}{\nobreak\hbox{-}}
\setlength{\emergencystretch}{2em}
\multlinegap0pt
\usepackage{bm}
\usepackage{bbm}
\usepackage{dsfont}
\usepackage{xspace}
\usepackage{cancel}
\usepackage{mathtools}
\usepackage{amsthm}
\makeatletter
\@ifundefined{thm}{\newtheorem{thm}{Thm}}{}
\@ifundefined{lemma}{\newtheorem{lemma}[thm]{Lemma}}{}
\makeatother
\DeclareMathOperator{\Crit}{Crit}
\DeclareMathOperator{\Spec}{Spec}
\newcommand{\R}{{\mathbb{R}}}
\newcommand{\Z}{{\mathbb{Z}}}
\newcommand{\cP}{{\mathcal{P}}}
\newcommand{\cQ}{{\mathcal{Q}}}
\newcommand{\fc}{{:\ }}
\title[On the Lagrangian Hofer geometry of Clifford tori]{On the Lagrangian Hofer geometry of Clifford tori}
\author{Frol Zapolsky}
\date{Source: arXiv:2508.02501v1 (CC BY 4.0)}
\begin{document}
\maketitle

\noindent\emph{Source and licence.} On the Lagrangian Hofer geometry of Clifford tori. Version arXiv:2508.02501v1 (CC BY 4.0), distributed under the Creative Commons Attribution 4.0 International licence, as recorded in the arXiv licence metadata for this exact version and shown on the abstract page https://arxiv.org/abs/2508.02501v1. Full rights notice: licenses/arxiv-2508.02501-CC-BY-4.0.txt.\par\smallskip

\noindent\textit{Editorial scope.} Counted unit: the Lemma whose source label is \texttt{lem:boundedness\_Lag\_sp\_invt} in the article (source0.tex lines 592--593, source label \texttt{lem:boundedness\_Lag\_sp\_invt}), delivered together with the article's own complete original proof of it (source0.tex lines 594--616, one \texttt{proof} environment of 23 lines). No mathematics is written, repaired or re-derived: statement and proof are the article's own text line for line. Required context retained uncounted: none. Every definition, display and label of those lines is retained unchanged. Omitted from this package and not redistributed: 1--591, 617--677. Nothing inside a retained excerpt is omitted or altered: every retained body is delivered exactly as the article's lines are written, so no omission receipt of any kind accompanies this package. Citations: the retained excerpts carry no citation command at all, so no bibliography entry of the article is transcribed and no reference is redistributed. Reference adaptation: none was needed and none was made; every reference inside the delivered unit resolves inside it, so no unresolved reference or citation is produced. Printed slips kept literal: no printed word, symbol, hypothesis or number of the counted statement or of its proof was altered. Layout and class: the article's own class is \texttt{article} with packages including xy, inputenc, amsfonts, amsmath, amssymb, amscd, amsthm, stmaryrd, indentfirst, geometry, faktor, epsfig; the wrapper is the lane's pinned \texttt{amsart} editorial wrapper at 10pt on A4, which restates the theorem-like environments the retained text needs and adds the article's own macro definitions that the retained text needs (\texttt{\textbackslash Crit}, \texttt{\textbackslash R}, \texttt{\textbackslash Spec}, \texttt{\textbackslash Z}, \texttt{\textbackslash cP}, \texttt{\textbackslash cQ}, \texttt{\textbackslash fc}), with extra wrapper packages (bm, bbm, dsfont, xspace, cancel, mathtools). The delivered numbering is the unit's own. Assets: no retained span contains a figure environment, a graphics-inclusion command, a PDF-page inclusion, a table or a TikZ picture; no image file is copied into this package, the package declares no asset dependency and no image file is redistributed with it. Licence: version arXiv:2508.02501v1 of this article is distributed under the Creative Commons Attribution 4.0 International licence, as recorded in the arXiv licence metadata for this exact version and shown on the abstract page https://arxiv.org/abs/2508.02501v1. A full rights notice is delivered at licenses/arxiv-2508.02501-CC-BY-4.0.txt. Characters: every retained excerpt is pure ASCII, so the delivered engine, XeLaTeX with its default Latin Modern text font, reports no missing character.
\begin{lemma}\label{lem:boundedness_Lag_sp_invt}Let $(M,\omega)$ be a closed connected symplectic manifold and let $L$ be a Lagrangian with monotonicity constant $\tau$. Let $R,R'$ be arbitrary ground rings, let $\cQ,\cQ'$ be local systems over those rings, and suppose $QH_*(L;\cQ)\neq 0\neq QH_*(L;\cQ')$. Then  for any Hamiltonian $G^0$ on $M$ with $\phi_{G^0}(L)=L$ we have $|\ell_L^{\cQ}(G^0)-\ell_L^{\cQ'}(G^0)|\leq \tau\cdot\dim L$ .
\end{lemma}

\begin{proof}
  Let $n=\dim L = \frac 12 \dim M$. Denote $\ell=\ell^\cQ_L$, $\ell'=\ell^{\cQ'}_L$. Let $G^0$ be a Hamiltonian such that $\phi_{G^0}(L) = L$, which we can assume satisfies $G^0_t\equiv 0$ for $t$ near $1$, by suitable time reparametrization. Let $a = \ell(G^0)$.

  Fix $\epsilon > 0$ and let $G$ be a perturbation of $G^0$ obtained as follows. Suppose $\delta > 0$ is such that $G^0_t\equiv 0$ for $t\in[1-\delta,1]$. Let $f \in C^\infty(L)$ be a Morse function, and let $H \in C^\infty(M)$ be obtained by suitably cutting off $\pi^*f$, where $\pi$ is the bundle projection on a Weinstein neighborhood of $L$, so that locally in that neighborhood, $\phi^t_H(L) = \phi_{\pi^*f}^t(L)$ for $t \in [0,1]$. Now time-reparametrize $H$ to obtain a Hamiltonian $H'$ satisfying $H'_t\equiv 0$ for $t \leq 1-\delta$ and replace the zero portion of $G^0_t$ for $t \in [1-\delta,1]$ by $H'$. Let us denote the resulting Hamiltonian by $G$. Suitably rescaling $f$, we can achieve $\|H'\|_{C^0} < \epsilon$, and thus $\|G^0-G\|_{C^0} < \epsilon$.

  We have $|\ell(G) - a| = |\ell(G) - \ell(G^0)| < \epsilon$ by continuity. By spectrality, $a \in \Spec(G^0:L)$, whence
  $$\Spec(G^0:L) = a + \mathsf A\cdot \Z\,,$$
  where $\mathsf A >0$ is the positive generator of the subgroup $\langle \omega,\pi_2(M,L)\rangle \subset \R$. Thanks to our construction of the perturbation $G$, the set $\cP_L(G)$ consists exactly of those curves $\gamma \fc [0,1] \to M$ which satisfy $\gamma(t) = \phi_{G^0}^t(\gamma(0))$ and $\gamma(1) \in \Crit f$. A straightforward action calculation then implies that
  $$\Spec(G:L) \subset (a-\epsilon,a+\epsilon) + \mathsf A \cdot \Z\,.$$
  Let us refer to the capped orbits of $G$ whose actions are contained in an interval $(b,c)$ as orbits of $G$ in the window $(b,c)$. By construction, since $G$ is nondegenerate in the sense that $\phi_G(L)$ intersects $L$ transversely, there exists $i \in \Z$ such that the Viterbo--Maslov indices of the orbits of $G$ in the window $(a-\epsilon,a+\epsilon)$ all lie in the set
  $$\{i,i+1,\dots,i+n\}\,.$$
  By spectrality, $\ell(G)$, being in the interval $(a-\epsilon,a+\epsilon)$, is the action relative to $G$ of one of the orbits in this window, which moreover has index $n$, whence
  $$n \in \{i,\dots,i+n\}\,,$$
  which implies $i \in \{0,\dots,n\}$. Since $L$ is monotone, the orbits of $G$ in the window $(a-\epsilon,a+\epsilon)+k\mathsf A$, where $k \in \Z$, all have indices in the set
  $$\{i,\dots,i+n\} + kN_L\,.$$
  For $|k|>\frac{n}{N_L}$ we have
  $$n \notin \{i,\dots,i+n\} + kN_L\,.$$
  Again by spectrality, $\ell'(G)$ is the action of an orbit of $G$ of index $n$, and from the above it follows that it must lie in a window $(a-\epsilon,a+\epsilon)+k\mathsf A$ where $|k|\leq \frac{n}{N_L}$. It follows that
  $$|\ell'(G) - a| < \epsilon + \frac{n}{N_L}\cdot\mathsf A\,.$$
  Taking $\epsilon \to 0$, we have $\ell'(G) \to \ell'(G^0)$ and therefore
  $$|\ell(G^0) - \ell'(G^0)| \leq \frac{n\mathsf A}{N_L}=\tau\cdot\dim L\,,$$
  where we used the fact that $\tau=\frac{\mathsf A}{N_L}$. The lemma is proved.
\end{proof}

\end{document}
